\documentclass[11pt,a4paper]{article}

\usepackage[T1]{fontenc}
\usepackage[utf8]{inputenc}
\usepackage[english]{babel}
\usepackage{microtype}
\usepackage{geometry}
\geometry{margin=2.5cm}
\usepackage{amsmath, amssymb}
\usepackage{booktabs}
\usepackage{longtable}
\usepackage[table]{xcolor}
\usepackage{listings}
\usepackage{hyperref}
\hypersetup{
    colorlinks=true,
    urlcolor=blue,
    linkcolor=black,
    citecolor=blue,
}
\usepackage{graphicx}
\usepackage{caption}
\usepackage{titlesec}

\titleformat{\section}{\large\bfseries}{\thesection.}{0.6em}{}
\titleformat{\subsection}{\normalsize\bfseries}{\thesubsection.}{0.6em}{}

\definecolor{codebg}{rgb}{0.96,0.96,0.96}
\definecolor{codekw}{rgb}{0.10,0.30,0.70}
\definecolor{codestr}{rgb}{0.50,0.10,0.10}
\definecolor{codecmt}{rgb}{0.30,0.55,0.30}
\lstdefinestyle{py}{
    language=Python,
    basicstyle=\ttfamily\small,
    backgroundcolor=\color{codebg},
    keywordstyle=\color{codekw}\bfseries,
    stringstyle=\color{codestr},
    commentstyle=\color{codecmt}\itshape,
    breaklines=true,
    showstringspaces=false,
    columns=fullflexible,
    frame=single,
    framerule=0pt,
    framesep=4pt,
    xleftmargin=4pt,
}
\lstset{style=py}

\title{\textbf{PyRINEX 4.0 Manual}\\
\large A Python toolkit for batch processing and quality\\
checking of RINEX 2/3 GNSS data}
\author{Han Jinzhen \\ \href{mailto:hanjinzhen9@gmail.com}{hanjinzhen9@gmail.com}}
\date{\today}

\begin{document}
\maketitle
\tableofcontents
\bigskip

\section{Introduction}

PyRINEX is a Python package for batch processing and quality-checking
GNSS observation files in the RINEX 2.xx and 3.xx formats. It targets
two use cases that occur repeatedly in geodetic and surveying
workflows: (i) consolidating large, inconsistently-named archives of
RINEX files into a clean directory tree, and (ii) running standard
quality-check products---multipath, ionospheric drift, cycle-slip
detection, sky plot, signal plot---on every file in a folder without
having to drop into a GUI tool.

The package is released under the Apache License 2.0; the canonical
citation is the PeerJ Computer Science paper that introduced it
\cite{han2024pyrinex}.

This manual documents PyRINEX 4.0, which is a substantial rewrite of
the v3.x release. The high-level API and command surface are
backwards-compatible (legacy names continue to work and emit
\texttt{DeprecationWarning}), but every public function now returns
plain Python objects rather than JSON strings, the package layout has
been split into focused submodules, the typo'd module name
\texttt{QulityCheck} has been corrected to \texttt{quality\_check},
and a number of bugs identified post-publication have been fixed
(see \S\ref{sec:changelog}).

\section{Installation}
\label{sec:install}

PyRINEX is published on PyPI:

\begin{lstlisting}[language=bash]
pip install PyRINEX
\end{lstlisting}

To install the development version from source:

\begin{lstlisting}[language=bash]
git clone https://github.com/geumjin99/PyRINEX
cd PyRINEX
pip install -e .[dev]
\end{lstlisting}

Python~3.9 or newer is required. Runtime dependencies are
\texttt{numpy}, \texttt{matplotlib}, \texttt{seaborn},
\texttt{chardet}, and \texttt{python-dateutil}. The development
extras add \texttt{pytest} and \texttt{ruff} for tests and linting.

\section{Quick start}

After installation, import the modules and call them on a RINEX file:

\begin{lstlisting}
from PyRINEX import reader, data_management, quality_check

header = reader.read_obs_header("test0010.24o")
obs    = reader.read_obs("test0010.24o")
nav    = reader.read_nav("test0010.24n")

quality_check.quality_check("test0010.24o")
\end{lstlisting}

\noindent
The next sections describe each module in detail.

\section{Reader module}

The \texttt{PyRINEX.reader} module turns RINEX text files into native
Python data structures. Three high-level functions cover the typical
needs:

\begin{longtable}{p{0.32\linewidth}p{0.61\linewidth}}
\toprule
\textbf{Function} & \textbf{Description} \\
\midrule
\texttt{read\_obs\_header(path)} &
Parse the header section of an observation file; returns an
\texttt{ObsHeader} dataclass. \\
\addlinespace
\texttt{read\_obs(path)} &
Parse the body of an observation file; returns a dict mapping
epoch-label strings to per-PRN observation dicts. \\
\addlinespace
\texttt{read\_nav(path)} &
Parse a GPS broadcast-navigation file; returns a dict mapping
each PRN to a list of broadcast ephemeris records. \\
\addlinespace
\texttt{read\_nav\_iono(path)} &
Extract the eight Klobuchar ionosphere model parameters from a
GPS navigation file. \\
\bottomrule
\end{longtable}

\subsection{Header}

\texttt{ObsHeader} (defined in \texttt{PyRINEX.reader}) is a frozen
dataclass with the following fields:

\begin{longtable}{p{0.32\linewidth}p{0.61\linewidth}}
\toprule
\textbf{Field} & \textbf{Meaning} \\
\midrule
\texttt{version} & RINEX major version (\texttt{2} or \texttt{3}). \\
\texttt{type} & Observation system letter (\texttt{G}, \texttt{M}, \dots). \\
\texttt{MARKER\_NAME} & \texttt{[line, value]}; line index is preserved
so \texttt{clean\_rinex} can rewrite this header line in place. \\
\texttt{MARKER\_NUMBER} & As above. \\
\texttt{receiver\_type} / \texttt{antenna\_type} & As above. \\
\texttt{APPROX\_POSITION\_XYZ} & ECEF coordinates as a 3-tuple of
floats, in metres. \\
\texttt{TIME\_OF\_FIRST\_OBS} / \texttt{TIME\_OF\_LAST\_OBS} & Hyphen-joined
time strings; both are kept distinct (a bug in v3.x clobbered
one with the other). \\
\texttt{ObsTypes} & List of observation codes for v2; per-system
dict for v3. \\
\texttt{PRNS} & Sorted list of every PRN that appears in the body of
the file. \\
\texttt{END\_OF\_HEADER} & Line index of \texttt{END OF HEADER}. \\
\bottomrule
\end{longtable}

\subsection{Observations}

\texttt{read\_obs} returns a \texttt{dict} keyed by epoch label
(``\texttt{yy~mm~dd~hh~mi~ss.ss}''). Each value is itself a dict
containing the entry \texttt{sat\_num} and one entry per PRN whose
value is a dict from observation-type to the verbatim 16-character
RINEX field string. Values that are missing in the file are
substituted with the sentinel \texttt{"~~~~~~~~~0.000~~"}. Downstream
quality-check routines treat this sentinel as ``no observation''.

\subsection{Navigations}

\texttt{read\_nav} returns a dict from PRN to a list of broadcast
ephemerides. Each ephemeris is a dict whose keys are descriptive
parameter names (``Toe Time of Ephemeris'', ``sqrt(A)'',
``Cuc'', \dots) and whose values are the verbatim 19-character
fields straight from the navigation file. To convert a field to a
floating-point value, use the \texttt{\_to\_float} helper from
\texttt{PyRINEX.orbit} (which understands Fortran-style ``D''
exponents):

\begin{lstlisting}
from PyRINEX.orbit import _to_float
toe = _to_float(nav["G01"][0]["Toe Time of Ephemeris"])
\end{lstlisting}

\subsection{Legacy JSON API}

For backwards compatibility with code written against v3.x, the
following functions return JSON-encoded strings rather than dicts:
\texttt{oheader(path)}, \texttt{observations(path)},
\texttt{navigations(path)}, \texttt{nheader(path)}. They emit
\texttt{DeprecationWarning} on every call. New code should prefer the
``\texttt{read\_*}'' functions which return native objects directly
and avoid the JSON round-trip entirely.

\section{DataManagement module}

\texttt{PyRINEX.data\_management} provides bulk file-discovery and
RINEX-cleaning helpers.

\subsection{find\_rinex (formerly DataFinding)}
\label{sec:findrinex}

Given a root directory, a list of folder-name keywords and a file
extension, \texttt{find\_rinex} returns every matching path:

\begin{lstlisting}
from PyRINEX.data_management import find_rinex

paths = find_rinex(
    root_path="/data/2008",
    keywords=["Gyeong-gi", "Unified control points"],
    extension=".08o",
)
\end{lstlisting}

The match is case-insensitive on the extension; the keywords must all
be substrings of the folder path of each candidate file.

\subsection{clean\_rinex (formerly DataCleaning)}

\texttt{clean\_rinex} performs five corrections on a list of RINEX
observation files and copies the matching navigation file alongside:

\begin{enumerate}
    \item Truncate the marker name to four characters
    (right-padding with ``\texttt{a}'' if necessary).
    \item Substitute receiver-type strings according to a user-provided
    \texttt{ReceiverLibrary.csv}.
    \item Substitute antenna-type strings according to
    \texttt{AntennaLibrary.csv}.
    \item Optionally fill in \texttt{NONE} as the antenna radome.
    \item Replace non-ASCII characters in the first twenty header
    lines with ``\texttt{-}''.
\end{enumerate}

The library CSVs are simple two-column tables of the form
``\texttt{wrong?right}''---one entry per line, comma-free in either
column---which lets users add their own entries without touching
PyRINEX's source. A typical receiver library entry:

\begin{lstlisting}
TRIMBLENETR9 ?TRIMBLE NETR9
TPS GB1000   ?TPS GB-1000
\end{lstlisting}

After running, \texttt{clean\_rinex} writes a CSV report at
\texttt{\textit{output\_root}/report.csv} summarising every file
processed, its origin and target paths, and the changes that were
made.

\subsection{Coordinate conversion}

The WGS-84 ECEF$\leftrightarrow$BLH conversions used by
\texttt{clean\_rinex} are exposed independently in
\texttt{PyRINEX.coordinates}:

\begin{lstlisting}
from PyRINEX.coordinates import xyz_to_blh, blh_to_xyz

lon, lat, h = xyz_to_blh(-3173543.0, 4134173.3, 3666275.5,
                         radians=False)
\end{lstlisting}

\noindent
The forward conversion uses the closed-form Bowring iteration with a
0.1\,m convergence threshold on height.

\section{QualityCheck module}

\texttt{PyRINEX.quality\_check} provides every per-file quality-check
routine and the two batch drivers \texttt{quality\_check} and
\texttt{batch\_quality\_check}.

\begin{longtable}{p{0.34\linewidth}p{0.59\linewidth}}
\toprule
\textbf{Function} & \textbf{Description} \\
\midrule
\texttt{satellite\_signal\_plot(path)} & Per-PRN signal-presence
raster (Figure~\ref{fig:signal-plot}). \\
\addlinespace
\texttt{azimuth\_elevation(path)} & Azimuth and elevation of every GPS
satellite per epoch. Writes a sky-plot PNG and a CSV. \\
\addlinespace
\texttt{multipath\_iod(path)} & MP1 / MP2 multipath, ionospheric delay
(ION) and drift (IOD). Writes four PNGs and two CSVs. \\
\addlinespace
\texttt{cycle\_slip(path)} & Carrier-phase second-difference
cycle-slip indicator. Writes one PNG and one CSV. \\
\addlinespace
\texttt{quality\_check(path)} & Run all four routines on a single
file. \\
\addlinespace
\texttt{batch\_quality\_check(root, keywords, ext)} & Run
\texttt{quality\_check} on every file matching the
\texttt{find\_rinex} criteria of \S\ref{sec:findrinex}. \\
\bottomrule
\end{longtable}

\subsection{Multipath observable}

For each constellation supported by \texttt{frequency\_pair} (GPS,
GLONASS, Galileo, SBAS), the L1 and L2 multipath observables are
computed from the dual-frequency code--phase combination
\cite{estey1999teqc}:
\begin{align}
    \mathrm{MP_1}
    &= P_1 - \left(\frac{2}{\gamma - 1} + 1\right) \lambda_1 L_1
        + \frac{2}{\gamma - 1} \lambda_2 L_2, \\
    \mathrm{MP_2}
    &= P_2 - \frac{2\gamma}{\gamma - 1} \lambda_1 L_1
        + \left(\frac{2\gamma}{\gamma - 1} - 1\right) \lambda_2 L_2,
\end{align}
where $\gamma = (f_1 / f_2)^2$. Per continuous arc the per-arc mean
of MP1 and MP2 is removed, isolating the multipath component from the
arbitrary integer ambiguity bias.

\subsection{Geometry-free combination and ionospheric drift}

The geometry-free phase combination
\begin{equation}
    \mathrm{ION} = \frac{\gamma}{\gamma - 1}
                    (\lambda_1 L_1 - \lambda_2 L_2)
\end{equation}
is differenced between epochs to produce the
ionospheric drift indicator IOD = $d\mathrm{ION}/dt$.

\subsection{Cycle-slip detector}

PyRINEX uses the second-difference of the geometry-free combination
to flag cycle slips in carrier phase:
\begin{equation}
    \Delta^2 \Phi_{\mathrm{IF}}(t) = \Phi_{\mathrm{IF}}(t-1)
        - 2\Phi_{\mathrm{IF}}(t) + \Phi_{\mathrm{IF}}(t+1).
\end{equation}
A non-zero second difference outside numerical noise is taken to mean
either a slip or a code-range outlier; the latter is rejected by an
absolute-value threshold.

\subsection{GPS satellite position}

The azimuth/elevation product depends on the ECEF position of each
GPS satellite, computed from the broadcast ephemeris using the
algorithm in IS-GPS-200, Table~20-IV. PyRINEX exposes that
calculation as \texttt{PyRINEX.orbit.gps\_satellite\_position(eph,
t\_k)} so users can call it directly:

\begin{lstlisting}
from PyRINEX.reader import read_nav
from PyRINEX.orbit import gps_satellite_position

nav = read_nav("test0010.24n")
x, y, z = gps_satellite_position(nav["G01"][0], t_k=0.0)
\end{lstlisting}

\subsection{Output products}

When run on a file \texttt{stub.24o}, \texttt{quality\_check} writes
the following sibling files (matching the published manual screenshot):

\begin{itemize}
    \item \texttt{stubSignalPlot.jpg}
    \item \texttt{stubskyplot.png} (only when a matching nav file
      is found)
    \item \texttt{stubaziele.csv}
    \item \texttt{stubMP1\_plot.png} / \texttt{stubMP2\_plot.png}
    \item \texttt{stubION\_plot.png} / \texttt{stubIOD\_plot.png}
    \item \texttt{stubmp1mp2.csv} / \texttt{stubIonIod.csv}
    \item \texttt{stubCycleSlipCarrier.png} /
      \texttt{stubCScarrier.csv}
\end{itemize}

\section{Practical use cases}

\subsection{Bulk data cleaning}

Suppose you have a folder hierarchy
\texttt{2008/<region>/<point\_type>/<date>/RINEX/<station>/}
that contains observation files with marker names that are too long,
non-ASCII characters in the header, and inconsistent receiver-type
spellings. The following script consolidates the Gyeong-gi
``Unified control points'' subset into a tidy date-indexed tree:

\begin{lstlisting}
from PyRINEX.data_management import clean_rinex, find_rinex

paths = find_rinex(
    root_path="/data/2008",
    keywords=["Gyeong-gi", "Unified control points"],
    extension=".08o",
)
clean_rinex(
    rinex_files=paths,
    receiver_library_path="ReceiverLibrary.csv",
    antenna_library_path="AntennaLibrary.csv",
    output_root="/data/2008-clean",
    radome=True,
)
\end{lstlisting}

\subsection{Batch quality check}

Once a folder is clean, run quality-check on every file:

\begin{lstlisting}
from PyRINEX.quality_check import batch_quality_check

batch_quality_check(
    root_path="/data/2008-clean",
    keywords=[],
    extension=".08o",
)
\end{lstlisting}

\section{Changelog}
\label{sec:changelog}

\subsection*{4.0.0}

\begin{itemize}
    \item Module \texttt{QulityCheck} has been renamed to
      \texttt{quality\_check} (typo fix). The old name is preserved
      as a deprecation shim.
    \item Module \texttt{DataManagement} has been renamed to
      \texttt{data\_management} (PEP-8 naming). The old name is
      preserved as a deprecation shim.
    \item \texttt{read\_obs\_header} / \texttt{read\_obs} /
      \texttt{read\_nav} return native Python objects instead of
      JSON strings. The legacy \texttt{oheader} / \texttt{observations}
      / \texttt{navigations} helpers continue to return JSON for
      callers that wrap them with \texttt{json.loads}.
    \item Bug: \texttt{TIME OF LAST OBS} no longer overwrites
      \texttt{TIME OF FIRST OBS}.
    \item Bug: \texttt{abs(MP1 > threshold)} was an always-true
      \texttt{abs(bool)}; corrected to
      \texttt{abs(MP1) > threshold}.
    \item Bug: hard-coded ``\texttt{\textbackslash\textbackslash}''
      path separators in the cleaning routine have been replaced with
      \texttt{os.path.join} (now works on Linux and macOS).
    \item Bug: the inner-\texttt{return} indentation in
      \texttt{file\_name\_extesion\_judgement} that caused only the
      first directory to be scanned has been fixed.
    \item Bug: \texttt{new\_marker.replace(" ", "a")} was a no-op
      (strings are immutable); the result is now assigned back.
    \item Bug: aliasing between \texttt{prns} and \texttt{fieldnames}
      in the multipath/cycle-slip output writers has been removed.
    \item The four hand-unrolled per-constellation extraction blocks
      in \texttt{LandC} have been replaced by a config-driven loop.
    \item Distribution metadata has moved from \texttt{setup.py} to
      \texttt{pyproject.toml} (PEP~621). The \texttt{yourusername}
      placeholder URL has been corrected.
    \item Tests: a \texttt{pytest} suite with synthetic RINEX 2/3
      fixtures runs end-to-end without external data.
\end{itemize}

\bigskip
\bibliographystyle{plain}
\bibliography{manual}

\end{document}
